\documentclass[11pt,a4paper]{article}
% main (standalone preamble for editing)
% vim: nowrap: spell:
% ══════════════════════════════════════════════════════════════════════════════
% Speed Maths --- shared preamble
% Included by every sheet and answer file via:  % main (standalone preamble for editing)
% vim: nowrap: spell:
% ══════════════════════════════════════════════════════════════════════════════
% Speed Maths --- shared preamble
% Included by every sheet and answer file via:  % main (standalone preamble for editing)
% vim: nowrap: spell:
% ══════════════════════════════════════════════════════════════════════════════
% Speed Maths --- shared preamble
% Included by every sheet and answer file via:  \input{../../shared/preamble}
% Editing THIS file changes the look of every sheet/answer across every pillar.
% ══════════════════════════════════════════════════════════════════════════════

\usepackage[a4paper, top=25mm, bottom=25mm, left=22mm, right=22mm]{geometry}
\usepackage{amsmath, amssymb}
\usepackage{enumitem}
\usepackage{booktabs}
\usepackage{array}
\usepackage{parskip}
\usepackage{microtype}
\usepackage{fancyhdr}
\usepackage{titlesec}
\usepackage{mdframed}
\usepackage{xcolor}
\usepackage[hidelinks]{hyperref}
\usepackage{graphicx}
\usepackage{tikz}
\usepackage{pgfplots}
\pgfplotsset{compat=1.18}
\usetikzlibrary{calc, angles, quotes}
\tikzset{
  speedgeom/.style={line width=0.8pt, draw=darkgray, line cap=round, line join=round},
  speedaux/.style={line width=0.5pt, dashed, draw=medgray},
  speedpoint/.style={circle, fill=darkgray, inner sep=1.2pt}
}

% ── Colours ───────────────────────────────────────────────────────────────────
\definecolor{darkgray}{gray}{0.15}
\definecolor{medgray}{gray}{0.45}
\definecolor{lightgray}{gray}{0.93}
\definecolor{boxgray}{gray}{0.88}

% ── Section formatting ──────────────────────────────────────────────────────────
\titleformat{\section}[block]
  {\normalfont\large\bfseries}{}
  {0pt}{}[\vspace{2pt}\hrule\vspace{6pt}]
\titlespacing*{\section}{0pt}{14pt}{4pt}

% ── List formatting ─────────────────────────────────────────────────────────────
\setlist[enumerate,1]{label=\textbf{\arabic*.}, leftmargin=2em, itemsep=10pt}

% ── Branding constants (edit here, not per-file) ────────────────────────────────
\newcommand{\SpeedExamLine}{\textbf{Speed Maths:} Competition \& Exam Prep}
\newcommand{\SpeedCredit}{Series by \href{https://www.linkedin.com/in/cerealdev/}{David Akinyele-Aje}}

% ── Header / footer ──────────────────────────────────────────────────────────────
% Usage (once, after \input{../../shared/preamble} and before \begin{document}):
%   \SpeedHeader{Algebra}{3}
% First arg  = pillar name shown as "Speed <Pillar> --- Daily Drill"
% Second arg = worksheet number
\pagestyle{fancy}
\newcommand{\SpeedHeader}[2]{%
  \fancyhf{}%
  \renewcommand{\headrulewidth}{0.4pt}%
  \setlength{\headheight}{14pt}%
  \fancyhead[L]{\small\textbf{Speed #1 --- Daily Drill}}%
  \fancyhead[R]{\small Worksheet \##2}%
  \fancyfoot[L]{\small \SpeedExamLine}%
  \fancyfoot[C]{\small\thepage}%
  \fancyfoot[R]{\small \SpeedCredit}%
  \renewcommand{\footrulewidth}{0.4pt}%
}

% ── Title block (used at the top of every sheet AND answer file) ───────────────
% Usage: \SpeedTitleBlock{Daily Algebra Drill \#3}{Jane Doe}
% %% AGENTS: When generating a new sheet, ask the human contributor for the name they want credited in argument 2.
% For answer files, pass the "--- Answers \& Investigations" suffix in the arg.
\newcommand{\SpeedTitleBlock}[2]{%
  \begin{center}
    {\LARGE\bfseries #1}\\[4pt]
    {\normalsize\itshape Sheet by #2}\\[6pt]
    \hrule\vspace{4pt}
    \begin{tabular}{llcc}
      \toprule
      \textbf{Section} & \textbf{Topic} & \textbf{Questions} & \textbf{Target time}\\
      \midrule
      A & Rapid Recognition          & 10 & 2 min 30 sec \\
      B & Manipulation Drills        & 10 & 8 min        \\
      C & Substitution \& Structure  & 8  & 10 min       \\
      D & Challenge                  & 5  & 15 min       \\
      \bottomrule
    \end{tabular}
    \vspace{4pt}
    \hrule
  \end{center}
}

% ── Sheet metadata: topic + the day's new tools ────────────────────────────────
% Usage (immediately after \SpeedTitleBlock, sheets only):
%   \SpeedMeta{Expanding, factorising, and the identity toolkit}
%             {difference of squares, sum/product form, completing the square}
%
% Arg 1 = the sheet's topic in one line. The website lists this instead of
%         "Sheet 03", so write the thing a reader would actually search for.
% Arg 2 = comma-separated, at most five: the tools this sheet introduces that
%         earlier sheets in the pillar did not. These become the website's tags.
%         Leave out anything the reader already met — the tags are meant to say
%         what is new today, not to summarise the sheet.
%
% Both are parsed straight out of this file by tools/build_website.py, so the
% sheet stays the single source of truth and the site cannot drift from it.
%
% This renders NOTHING. It is a declaration for the build script, not a piece of
% the sheet's layout: adding it to a sheet must not change that sheet's PDF, so
% the 25 existing sheets can be annotated without a single one of them being
% reissued. If the toolkit should also appear on the page, that is a separate
% decision about the sheet design — make it deliberately, here, not by leaving
% this macro printing something nobody asked it to.
\newcommand{\SpeedMeta}[2]{}

% ── Answer / Method / Investigate-further boxes (answer files only) ────────────
\newmdenv[
  backgroundcolor=lightgray,
  linecolor=medgray,
  linewidth=0.5pt,
  leftmargin=0pt, rightmargin=0pt,
  innerleftmargin=8pt, innerrightmargin=8pt,
  innertopmargin=5pt, innerbottommargin=5pt,
  skipabove=4pt, skipbelow=4pt
]{ansbox}

\newmdenv[
  backgroundcolor=white,
  linecolor=darkgray,
  linewidth=0.8pt,
  leftline=true, rightline=false, topline=false, bottomline=false,
  leftmargin=0pt, rightmargin=0pt,
  innerleftmargin=8pt, innerrightmargin=6pt,
  innertopmargin=4pt, innerbottommargin=4pt,
  skipabove=4pt, skipbelow=6pt
]{invbox}

\newcommand{\ans}[1]{%
  \begin{ansbox}
    \textbf{Answer:} #1
  \end{ansbox}}

\newcommand{\method}[1]{%
  {\small\textbf{Method:} #1}\par}

\newcommand{\inv}[1]{%
  \begin{invbox}
    \textbf{\small Investigate further:} {\small #1}
  \end{invbox}}

% ── Misc helpers ────────────────────────────────────────────────────────────────
\newcommand{\qn}[1]{\textbf{#1.}\;}

% Mistake-linking: attach a Top 5 Patterns / Common Traps bullet to the question
% label(s) in this sheet where it actually shows up, e.g. \seealso{B6, D2}
\newcommand{\seealso}[1]{\hfill\textit{\footnotesize(see #1)}}

% ── Graph options macro ─────────────────────────────────────────────────────────
% Usage: \GraphOptions{tikz code for A}{tikz code for B}{tikz code for C}{tikz code for D}
\newcommand{\GraphOptions}[4]{%
  \begin{center}
    \begin{tabular}{cc}
      \textbf{A)} & \textbf{B)} \\
      #1 & #2 \\[10pt]
      \textbf{C)} & \textbf{D)} \\
      #3 & #4
    \end{tabular}
  \end{center}
}

% ── Closing motivational rule + cross-promo footer (sheets only) ────────────────
% Usage: \SpeedClosing{"Quote text."}
\newcommand{\SpeedClosing}[1]{%
  \vspace{20pt}
  \begin{center}
  \rule{0.6\textwidth}{0.4pt}\\[4pt]
  {\small\itshape #1}\\[4pt]
  \rule{0.6\textwidth}{0.4pt}
  \end{center}
  \begin{center}
  {\small Found a mistake or want to contribute a sheet?}\\
  {\small Visit the open-source project at \href{https://github.com/The-CerealDev/speed-maths}{github.com/The-CerealDev/speed-maths}}
  \end{center}
}

% Editing THIS file changes the look of every sheet/answer across every pillar.
% ══════════════════════════════════════════════════════════════════════════════

\usepackage[a4paper, top=25mm, bottom=25mm, left=22mm, right=22mm]{geometry}
\usepackage{amsmath, amssymb}
\usepackage{enumitem}
\usepackage{booktabs}
\usepackage{array}
\usepackage{parskip}
\usepackage{microtype}
\usepackage{fancyhdr}
\usepackage{titlesec}
\usepackage{mdframed}
\usepackage{xcolor}
\usepackage[hidelinks]{hyperref}
\usepackage{graphicx}
\usepackage{tikz}
\usepackage{pgfplots}
\pgfplotsset{compat=1.18}
\usetikzlibrary{calc, angles, quotes}
\tikzset{
  speedgeom/.style={line width=0.8pt, draw=darkgray, line cap=round, line join=round},
  speedaux/.style={line width=0.5pt, dashed, draw=medgray},
  speedpoint/.style={circle, fill=darkgray, inner sep=1.2pt}
}

% ── Colours ───────────────────────────────────────────────────────────────────
\definecolor{darkgray}{gray}{0.15}
\definecolor{medgray}{gray}{0.45}
\definecolor{lightgray}{gray}{0.93}
\definecolor{boxgray}{gray}{0.88}

% ── Section formatting ──────────────────────────────────────────────────────────
\titleformat{\section}[block]
  {\normalfont\large\bfseries}{}
  {0pt}{}[\vspace{2pt}\hrule\vspace{6pt}]
\titlespacing*{\section}{0pt}{14pt}{4pt}

% ── List formatting ─────────────────────────────────────────────────────────────
\setlist[enumerate,1]{label=\textbf{\arabic*.}, leftmargin=2em, itemsep=10pt}

% ── Branding constants (edit here, not per-file) ────────────────────────────────
\newcommand{\SpeedExamLine}{\textbf{Speed Maths:} Competition \& Exam Prep}
\newcommand{\SpeedCredit}{Series by \href{https://www.linkedin.com/in/cerealdev/}{David Akinyele-Aje}}

% ── Header / footer ──────────────────────────────────────────────────────────────
% Usage (once, after % main (standalone preamble for editing)
% vim: nowrap: spell:
% ══════════════════════════════════════════════════════════════════════════════
% Speed Maths --- shared preamble
% Included by every sheet and answer file via:  \input{../../shared/preamble}
% Editing THIS file changes the look of every sheet/answer across every pillar.
% ══════════════════════════════════════════════════════════════════════════════

\usepackage[a4paper, top=25mm, bottom=25mm, left=22mm, right=22mm]{geometry}
\usepackage{amsmath, amssymb}
\usepackage{enumitem}
\usepackage{booktabs}
\usepackage{array}
\usepackage{parskip}
\usepackage{microtype}
\usepackage{fancyhdr}
\usepackage{titlesec}
\usepackage{mdframed}
\usepackage{xcolor}
\usepackage[hidelinks]{hyperref}
\usepackage{graphicx}
\usepackage{tikz}
\usepackage{pgfplots}
\pgfplotsset{compat=1.18}
\usetikzlibrary{calc, angles, quotes}
\tikzset{
  speedgeom/.style={line width=0.8pt, draw=darkgray, line cap=round, line join=round},
  speedaux/.style={line width=0.5pt, dashed, draw=medgray},
  speedpoint/.style={circle, fill=darkgray, inner sep=1.2pt}
}

% ── Colours ───────────────────────────────────────────────────────────────────
\definecolor{darkgray}{gray}{0.15}
\definecolor{medgray}{gray}{0.45}
\definecolor{lightgray}{gray}{0.93}
\definecolor{boxgray}{gray}{0.88}

% ── Section formatting ──────────────────────────────────────────────────────────
\titleformat{\section}[block]
  {\normalfont\large\bfseries}{}
  {0pt}{}[\vspace{2pt}\hrule\vspace{6pt}]
\titlespacing*{\section}{0pt}{14pt}{4pt}

% ── List formatting ─────────────────────────────────────────────────────────────
\setlist[enumerate,1]{label=\textbf{\arabic*.}, leftmargin=2em, itemsep=10pt}

% ── Branding constants (edit here, not per-file) ────────────────────────────────
\newcommand{\SpeedExamLine}{\textbf{Speed Maths:} Competition \& Exam Prep}
\newcommand{\SpeedCredit}{Series by \href{https://www.linkedin.com/in/cerealdev/}{David Akinyele-Aje}}

% ── Header / footer ──────────────────────────────────────────────────────────────
% Usage (once, after \input{../../shared/preamble} and before \begin{document}):
%   \SpeedHeader{Algebra}{3}
% First arg  = pillar name shown as "Speed <Pillar> --- Daily Drill"
% Second arg = worksheet number
\pagestyle{fancy}
\newcommand{\SpeedHeader}[2]{%
  \fancyhf{}%
  \renewcommand{\headrulewidth}{0.4pt}%
  \setlength{\headheight}{14pt}%
  \fancyhead[L]{\small\textbf{Speed #1 --- Daily Drill}}%
  \fancyhead[R]{\small Worksheet \##2}%
  \fancyfoot[L]{\small \SpeedExamLine}%
  \fancyfoot[C]{\small\thepage}%
  \fancyfoot[R]{\small \SpeedCredit}%
  \renewcommand{\footrulewidth}{0.4pt}%
}

% ── Title block (used at the top of every sheet AND answer file) ───────────────
% Usage: \SpeedTitleBlock{Daily Algebra Drill \#3}{Jane Doe}
% %% AGENTS: When generating a new sheet, ask the human contributor for the name they want credited in argument 2.
% For answer files, pass the "--- Answers \& Investigations" suffix in the arg.
\newcommand{\SpeedTitleBlock}[2]{%
  \begin{center}
    {\LARGE\bfseries #1}\\[4pt]
    {\normalsize\itshape Sheet by #2}\\[6pt]
    \hrule\vspace{4pt}
    \begin{tabular}{llcc}
      \toprule
      \textbf{Section} & \textbf{Topic} & \textbf{Questions} & \textbf{Target time}\\
      \midrule
      A & Rapid Recognition          & 10 & 2 min 30 sec \\
      B & Manipulation Drills        & 10 & 8 min        \\
      C & Substitution \& Structure  & 8  & 10 min       \\
      D & Challenge                  & 5  & 15 min       \\
      \bottomrule
    \end{tabular}
    \vspace{4pt}
    \hrule
  \end{center}
}

% ── Sheet metadata: topic + the day's new tools ────────────────────────────────
% Usage (immediately after \SpeedTitleBlock, sheets only):
%   \SpeedMeta{Expanding, factorising, and the identity toolkit}
%             {difference of squares, sum/product form, completing the square}
%
% Arg 1 = the sheet's topic in one line. The website lists this instead of
%         "Sheet 03", so write the thing a reader would actually search for.
% Arg 2 = comma-separated, at most five: the tools this sheet introduces that
%         earlier sheets in the pillar did not. These become the website's tags.
%         Leave out anything the reader already met — the tags are meant to say
%         what is new today, not to summarise the sheet.
%
% Both are parsed straight out of this file by tools/build_website.py, so the
% sheet stays the single source of truth and the site cannot drift from it.
%
% This renders NOTHING. It is a declaration for the build script, not a piece of
% the sheet's layout: adding it to a sheet must not change that sheet's PDF, so
% the 25 existing sheets can be annotated without a single one of them being
% reissued. If the toolkit should also appear on the page, that is a separate
% decision about the sheet design — make it deliberately, here, not by leaving
% this macro printing something nobody asked it to.
\newcommand{\SpeedMeta}[2]{}

% ── Answer / Method / Investigate-further boxes (answer files only) ────────────
\newmdenv[
  backgroundcolor=lightgray,
  linecolor=medgray,
  linewidth=0.5pt,
  leftmargin=0pt, rightmargin=0pt,
  innerleftmargin=8pt, innerrightmargin=8pt,
  innertopmargin=5pt, innerbottommargin=5pt,
  skipabove=4pt, skipbelow=4pt
]{ansbox}

\newmdenv[
  backgroundcolor=white,
  linecolor=darkgray,
  linewidth=0.8pt,
  leftline=true, rightline=false, topline=false, bottomline=false,
  leftmargin=0pt, rightmargin=0pt,
  innerleftmargin=8pt, innerrightmargin=6pt,
  innertopmargin=4pt, innerbottommargin=4pt,
  skipabove=4pt, skipbelow=6pt
]{invbox}

\newcommand{\ans}[1]{%
  \begin{ansbox}
    \textbf{Answer:} #1
  \end{ansbox}}

\newcommand{\method}[1]{%
  {\small\textbf{Method:} #1}\par}

\newcommand{\inv}[1]{%
  \begin{invbox}
    \textbf{\small Investigate further:} {\small #1}
  \end{invbox}}

% ── Misc helpers ────────────────────────────────────────────────────────────────
\newcommand{\qn}[1]{\textbf{#1.}\;}

% Mistake-linking: attach a Top 5 Patterns / Common Traps bullet to the question
% label(s) in this sheet where it actually shows up, e.g. \seealso{B6, D2}
\newcommand{\seealso}[1]{\hfill\textit{\footnotesize(see #1)}}

% ── Graph options macro ─────────────────────────────────────────────────────────
% Usage: \GraphOptions{tikz code for A}{tikz code for B}{tikz code for C}{tikz code for D}
\newcommand{\GraphOptions}[4]{%
  \begin{center}
    \begin{tabular}{cc}
      \textbf{A)} & \textbf{B)} \\
      #1 & #2 \\[10pt]
      \textbf{C)} & \textbf{D)} \\
      #3 & #4
    \end{tabular}
  \end{center}
}

% ── Closing motivational rule + cross-promo footer (sheets only) ────────────────
% Usage: \SpeedClosing{"Quote text."}
\newcommand{\SpeedClosing}[1]{%
  \vspace{20pt}
  \begin{center}
  \rule{0.6\textwidth}{0.4pt}\\[4pt]
  {\small\itshape #1}\\[4pt]
  \rule{0.6\textwidth}{0.4pt}
  \end{center}
  \begin{center}
  {\small Found a mistake or want to contribute a sheet?}\\
  {\small Visit the open-source project at \href{https://github.com/The-CerealDev/speed-maths}{github.com/The-CerealDev/speed-maths}}
  \end{center}
}
 and before \begin{document}):
%   \SpeedHeader{Algebra}{3}
% First arg  = pillar name shown as "Speed <Pillar> --- Daily Drill"
% Second arg = worksheet number
\pagestyle{fancy}
\newcommand{\SpeedHeader}[2]{%
  \fancyhf{}%
  \renewcommand{\headrulewidth}{0.4pt}%
  \setlength{\headheight}{14pt}%
  \fancyhead[L]{\small\textbf{Speed #1 --- Daily Drill}}%
  \fancyhead[R]{\small Worksheet \##2}%
  \fancyfoot[L]{\small \SpeedExamLine}%
  \fancyfoot[C]{\small\thepage}%
  \fancyfoot[R]{\small \SpeedCredit}%
  \renewcommand{\footrulewidth}{0.4pt}%
}

% ── Title block (used at the top of every sheet AND answer file) ───────────────
% Usage: \SpeedTitleBlock{Daily Algebra Drill \#3}{Jane Doe}
% %% AGENTS: When generating a new sheet, ask the human contributor for the name they want credited in argument 2.
% For answer files, pass the "--- Answers \& Investigations" suffix in the arg.
\newcommand{\SpeedTitleBlock}[2]{%
  \begin{center}
    {\LARGE\bfseries #1}\\[4pt]
    {\normalsize\itshape Sheet by #2}\\[6pt]
    \hrule\vspace{4pt}
    \begin{tabular}{llcc}
      \toprule
      \textbf{Section} & \textbf{Topic} & \textbf{Questions} & \textbf{Target time}\\
      \midrule
      A & Rapid Recognition          & 10 & 2 min 30 sec \\
      B & Manipulation Drills        & 10 & 8 min        \\
      C & Substitution \& Structure  & 8  & 10 min       \\
      D & Challenge                  & 5  & 15 min       \\
      \bottomrule
    \end{tabular}
    \vspace{4pt}
    \hrule
  \end{center}
}

% ── Sheet metadata: topic + the day's new tools ────────────────────────────────
% Usage (immediately after \SpeedTitleBlock, sheets only):
%   \SpeedMeta{Expanding, factorising, and the identity toolkit}
%             {difference of squares, sum/product form, completing the square}
%
% Arg 1 = the sheet's topic in one line. The website lists this instead of
%         "Sheet 03", so write the thing a reader would actually search for.
% Arg 2 = comma-separated, at most five: the tools this sheet introduces that
%         earlier sheets in the pillar did not. These become the website's tags.
%         Leave out anything the reader already met — the tags are meant to say
%         what is new today, not to summarise the sheet.
%
% Both are parsed straight out of this file by tools/build_website.py, so the
% sheet stays the single source of truth and the site cannot drift from it.
%
% This renders NOTHING. It is a declaration for the build script, not a piece of
% the sheet's layout: adding it to a sheet must not change that sheet's PDF, so
% the 25 existing sheets can be annotated without a single one of them being
% reissued. If the toolkit should also appear on the page, that is a separate
% decision about the sheet design — make it deliberately, here, not by leaving
% this macro printing something nobody asked it to.
\newcommand{\SpeedMeta}[2]{}

% ── Answer / Method / Investigate-further boxes (answer files only) ────────────
\newmdenv[
  backgroundcolor=lightgray,
  linecolor=medgray,
  linewidth=0.5pt,
  leftmargin=0pt, rightmargin=0pt,
  innerleftmargin=8pt, innerrightmargin=8pt,
  innertopmargin=5pt, innerbottommargin=5pt,
  skipabove=4pt, skipbelow=4pt
]{ansbox}

\newmdenv[
  backgroundcolor=white,
  linecolor=darkgray,
  linewidth=0.8pt,
  leftline=true, rightline=false, topline=false, bottomline=false,
  leftmargin=0pt, rightmargin=0pt,
  innerleftmargin=8pt, innerrightmargin=6pt,
  innertopmargin=4pt, innerbottommargin=4pt,
  skipabove=4pt, skipbelow=6pt
]{invbox}

\newcommand{\ans}[1]{%
  \begin{ansbox}
    \textbf{Answer:} #1
  \end{ansbox}}

\newcommand{\method}[1]{%
  {\small\textbf{Method:} #1}\par}

\newcommand{\inv}[1]{%
  \begin{invbox}
    \textbf{\small Investigate further:} {\small #1}
  \end{invbox}}

% ── Misc helpers ────────────────────────────────────────────────────────────────
\newcommand{\qn}[1]{\textbf{#1.}\;}

% Mistake-linking: attach a Top 5 Patterns / Common Traps bullet to the question
% label(s) in this sheet where it actually shows up, e.g. \seealso{B6, D2}
\newcommand{\seealso}[1]{\hfill\textit{\footnotesize(see #1)}}

% ── Graph options macro ─────────────────────────────────────────────────────────
% Usage: \GraphOptions{tikz code for A}{tikz code for B}{tikz code for C}{tikz code for D}
\newcommand{\GraphOptions}[4]{%
  \begin{center}
    \begin{tabular}{cc}
      \textbf{A)} & \textbf{B)} \\
      #1 & #2 \\[10pt]
      \textbf{C)} & \textbf{D)} \\
      #3 & #4
    \end{tabular}
  \end{center}
}

% ── Closing motivational rule + cross-promo footer (sheets only) ────────────────
% Usage: \SpeedClosing{"Quote text."}
\newcommand{\SpeedClosing}[1]{%
  \vspace{20pt}
  \begin{center}
  \rule{0.6\textwidth}{0.4pt}\\[4pt]
  {\small\itshape #1}\\[4pt]
  \rule{0.6\textwidth}{0.4pt}
  \end{center}
  \begin{center}
  {\small Found a mistake or want to contribute a sheet?}\\
  {\small Visit the open-source project at \href{https://github.com/The-CerealDev/speed-maths}{github.com/The-CerealDev/speed-maths}}
  \end{center}
}

% Editing THIS file changes the look of every sheet/answer across every pillar.
% ══════════════════════════════════════════════════════════════════════════════

\usepackage[a4paper, top=25mm, bottom=25mm, left=22mm, right=22mm]{geometry}
\usepackage{amsmath, amssymb}
\usepackage{enumitem}
\usepackage{booktabs}
\usepackage{array}
\usepackage{parskip}
\usepackage{microtype}
\usepackage{fancyhdr}
\usepackage{titlesec}
\usepackage{mdframed}
\usepackage{xcolor}
\usepackage[hidelinks]{hyperref}
\usepackage{graphicx}
\usepackage{tikz}
\usepackage{pgfplots}
\pgfplotsset{compat=1.18}
\usetikzlibrary{calc, angles, quotes}
\tikzset{
  speedgeom/.style={line width=0.8pt, draw=darkgray, line cap=round, line join=round},
  speedaux/.style={line width=0.5pt, dashed, draw=medgray},
  speedpoint/.style={circle, fill=darkgray, inner sep=1.2pt}
}

% ── Colours ───────────────────────────────────────────────────────────────────
\definecolor{darkgray}{gray}{0.15}
\definecolor{medgray}{gray}{0.45}
\definecolor{lightgray}{gray}{0.93}
\definecolor{boxgray}{gray}{0.88}

% ── Section formatting ──────────────────────────────────────────────────────────
\titleformat{\section}[block]
  {\normalfont\large\bfseries}{}
  {0pt}{}[\vspace{2pt}\hrule\vspace{6pt}]
\titlespacing*{\section}{0pt}{14pt}{4pt}

% ── List formatting ─────────────────────────────────────────────────────────────
\setlist[enumerate,1]{label=\textbf{\arabic*.}, leftmargin=2em, itemsep=10pt}

% ── Branding constants (edit here, not per-file) ────────────────────────────────
\newcommand{\SpeedExamLine}{\textbf{Speed Maths:} Competition \& Exam Prep}
\newcommand{\SpeedCredit}{Series by \href{https://www.linkedin.com/in/cerealdev/}{David Akinyele-Aje}}

% ── Header / footer ──────────────────────────────────────────────────────────────
% Usage (once, after % main (standalone preamble for editing)
% vim: nowrap: spell:
% ══════════════════════════════════════════════════════════════════════════════
% Speed Maths --- shared preamble
% Included by every sheet and answer file via:  % main (standalone preamble for editing)
% vim: nowrap: spell:
% ══════════════════════════════════════════════════════════════════════════════
% Speed Maths --- shared preamble
% Included by every sheet and answer file via:  \input{../../shared/preamble}
% Editing THIS file changes the look of every sheet/answer across every pillar.
% ══════════════════════════════════════════════════════════════════════════════

\usepackage[a4paper, top=25mm, bottom=25mm, left=22mm, right=22mm]{geometry}
\usepackage{amsmath, amssymb}
\usepackage{enumitem}
\usepackage{booktabs}
\usepackage{array}
\usepackage{parskip}
\usepackage{microtype}
\usepackage{fancyhdr}
\usepackage{titlesec}
\usepackage{mdframed}
\usepackage{xcolor}
\usepackage[hidelinks]{hyperref}
\usepackage{graphicx}
\usepackage{tikz}
\usepackage{pgfplots}
\pgfplotsset{compat=1.18}
\usetikzlibrary{calc, angles, quotes}
\tikzset{
  speedgeom/.style={line width=0.8pt, draw=darkgray, line cap=round, line join=round},
  speedaux/.style={line width=0.5pt, dashed, draw=medgray},
  speedpoint/.style={circle, fill=darkgray, inner sep=1.2pt}
}

% ── Colours ───────────────────────────────────────────────────────────────────
\definecolor{darkgray}{gray}{0.15}
\definecolor{medgray}{gray}{0.45}
\definecolor{lightgray}{gray}{0.93}
\definecolor{boxgray}{gray}{0.88}

% ── Section formatting ──────────────────────────────────────────────────────────
\titleformat{\section}[block]
  {\normalfont\large\bfseries}{}
  {0pt}{}[\vspace{2pt}\hrule\vspace{6pt}]
\titlespacing*{\section}{0pt}{14pt}{4pt}

% ── List formatting ─────────────────────────────────────────────────────────────
\setlist[enumerate,1]{label=\textbf{\arabic*.}, leftmargin=2em, itemsep=10pt}

% ── Branding constants (edit here, not per-file) ────────────────────────────────
\newcommand{\SpeedExamLine}{\textbf{Speed Maths:} Competition \& Exam Prep}
\newcommand{\SpeedCredit}{Series by \href{https://www.linkedin.com/in/cerealdev/}{David Akinyele-Aje}}

% ── Header / footer ──────────────────────────────────────────────────────────────
% Usage (once, after \input{../../shared/preamble} and before \begin{document}):
%   \SpeedHeader{Algebra}{3}
% First arg  = pillar name shown as "Speed <Pillar> --- Daily Drill"
% Second arg = worksheet number
\pagestyle{fancy}
\newcommand{\SpeedHeader}[2]{%
  \fancyhf{}%
  \renewcommand{\headrulewidth}{0.4pt}%
  \setlength{\headheight}{14pt}%
  \fancyhead[L]{\small\textbf{Speed #1 --- Daily Drill}}%
  \fancyhead[R]{\small Worksheet \##2}%
  \fancyfoot[L]{\small \SpeedExamLine}%
  \fancyfoot[C]{\small\thepage}%
  \fancyfoot[R]{\small \SpeedCredit}%
  \renewcommand{\footrulewidth}{0.4pt}%
}

% ── Title block (used at the top of every sheet AND answer file) ───────────────
% Usage: \SpeedTitleBlock{Daily Algebra Drill \#3}{Jane Doe}
% %% AGENTS: When generating a new sheet, ask the human contributor for the name they want credited in argument 2.
% For answer files, pass the "--- Answers \& Investigations" suffix in the arg.
\newcommand{\SpeedTitleBlock}[2]{%
  \begin{center}
    {\LARGE\bfseries #1}\\[4pt]
    {\normalsize\itshape Sheet by #2}\\[6pt]
    \hrule\vspace{4pt}
    \begin{tabular}{llcc}
      \toprule
      \textbf{Section} & \textbf{Topic} & \textbf{Questions} & \textbf{Target time}\\
      \midrule
      A & Rapid Recognition          & 10 & 2 min 30 sec \\
      B & Manipulation Drills        & 10 & 8 min        \\
      C & Substitution \& Structure  & 8  & 10 min       \\
      D & Challenge                  & 5  & 15 min       \\
      \bottomrule
    \end{tabular}
    \vspace{4pt}
    \hrule
  \end{center}
}

% ── Sheet metadata: topic + the day's new tools ────────────────────────────────
% Usage (immediately after \SpeedTitleBlock, sheets only):
%   \SpeedMeta{Expanding, factorising, and the identity toolkit}
%             {difference of squares, sum/product form, completing the square}
%
% Arg 1 = the sheet's topic in one line. The website lists this instead of
%         "Sheet 03", so write the thing a reader would actually search for.
% Arg 2 = comma-separated, at most five: the tools this sheet introduces that
%         earlier sheets in the pillar did not. These become the website's tags.
%         Leave out anything the reader already met — the tags are meant to say
%         what is new today, not to summarise the sheet.
%
% Both are parsed straight out of this file by tools/build_website.py, so the
% sheet stays the single source of truth and the site cannot drift from it.
%
% This renders NOTHING. It is a declaration for the build script, not a piece of
% the sheet's layout: adding it to a sheet must not change that sheet's PDF, so
% the 25 existing sheets can be annotated without a single one of them being
% reissued. If the toolkit should also appear on the page, that is a separate
% decision about the sheet design — make it deliberately, here, not by leaving
% this macro printing something nobody asked it to.
\newcommand{\SpeedMeta}[2]{}

% ── Answer / Method / Investigate-further boxes (answer files only) ────────────
\newmdenv[
  backgroundcolor=lightgray,
  linecolor=medgray,
  linewidth=0.5pt,
  leftmargin=0pt, rightmargin=0pt,
  innerleftmargin=8pt, innerrightmargin=8pt,
  innertopmargin=5pt, innerbottommargin=5pt,
  skipabove=4pt, skipbelow=4pt
]{ansbox}

\newmdenv[
  backgroundcolor=white,
  linecolor=darkgray,
  linewidth=0.8pt,
  leftline=true, rightline=false, topline=false, bottomline=false,
  leftmargin=0pt, rightmargin=0pt,
  innerleftmargin=8pt, innerrightmargin=6pt,
  innertopmargin=4pt, innerbottommargin=4pt,
  skipabove=4pt, skipbelow=6pt
]{invbox}

\newcommand{\ans}[1]{%
  \begin{ansbox}
    \textbf{Answer:} #1
  \end{ansbox}}

\newcommand{\method}[1]{%
  {\small\textbf{Method:} #1}\par}

\newcommand{\inv}[1]{%
  \begin{invbox}
    \textbf{\small Investigate further:} {\small #1}
  \end{invbox}}

% ── Misc helpers ────────────────────────────────────────────────────────────────
\newcommand{\qn}[1]{\textbf{#1.}\;}

% Mistake-linking: attach a Top 5 Patterns / Common Traps bullet to the question
% label(s) in this sheet where it actually shows up, e.g. \seealso{B6, D2}
\newcommand{\seealso}[1]{\hfill\textit{\footnotesize(see #1)}}

% ── Graph options macro ─────────────────────────────────────────────────────────
% Usage: \GraphOptions{tikz code for A}{tikz code for B}{tikz code for C}{tikz code for D}
\newcommand{\GraphOptions}[4]{%
  \begin{center}
    \begin{tabular}{cc}
      \textbf{A)} & \textbf{B)} \\
      #1 & #2 \\[10pt]
      \textbf{C)} & \textbf{D)} \\
      #3 & #4
    \end{tabular}
  \end{center}
}

% ── Closing motivational rule + cross-promo footer (sheets only) ────────────────
% Usage: \SpeedClosing{"Quote text."}
\newcommand{\SpeedClosing}[1]{%
  \vspace{20pt}
  \begin{center}
  \rule{0.6\textwidth}{0.4pt}\\[4pt]
  {\small\itshape #1}\\[4pt]
  \rule{0.6\textwidth}{0.4pt}
  \end{center}
  \begin{center}
  {\small Found a mistake or want to contribute a sheet?}\\
  {\small Visit the open-source project at \href{https://github.com/The-CerealDev/speed-maths}{github.com/The-CerealDev/speed-maths}}
  \end{center}
}

% Editing THIS file changes the look of every sheet/answer across every pillar.
% ══════════════════════════════════════════════════════════════════════════════

\usepackage[a4paper, top=25mm, bottom=25mm, left=22mm, right=22mm]{geometry}
\usepackage{amsmath, amssymb}
\usepackage{enumitem}
\usepackage{booktabs}
\usepackage{array}
\usepackage{parskip}
\usepackage{microtype}
\usepackage{fancyhdr}
\usepackage{titlesec}
\usepackage{mdframed}
\usepackage{xcolor}
\usepackage[hidelinks]{hyperref}
\usepackage{graphicx}
\usepackage{tikz}
\usepackage{pgfplots}
\pgfplotsset{compat=1.18}
\usetikzlibrary{calc, angles, quotes}
\tikzset{
  speedgeom/.style={line width=0.8pt, draw=darkgray, line cap=round, line join=round},
  speedaux/.style={line width=0.5pt, dashed, draw=medgray},
  speedpoint/.style={circle, fill=darkgray, inner sep=1.2pt}
}

% ── Colours ───────────────────────────────────────────────────────────────────
\definecolor{darkgray}{gray}{0.15}
\definecolor{medgray}{gray}{0.45}
\definecolor{lightgray}{gray}{0.93}
\definecolor{boxgray}{gray}{0.88}

% ── Section formatting ──────────────────────────────────────────────────────────
\titleformat{\section}[block]
  {\normalfont\large\bfseries}{}
  {0pt}{}[\vspace{2pt}\hrule\vspace{6pt}]
\titlespacing*{\section}{0pt}{14pt}{4pt}

% ── List formatting ─────────────────────────────────────────────────────────────
\setlist[enumerate,1]{label=\textbf{\arabic*.}, leftmargin=2em, itemsep=10pt}

% ── Branding constants (edit here, not per-file) ────────────────────────────────
\newcommand{\SpeedExamLine}{\textbf{Speed Maths:} Competition \& Exam Prep}
\newcommand{\SpeedCredit}{Series by \href{https://www.linkedin.com/in/cerealdev/}{David Akinyele-Aje}}

% ── Header / footer ──────────────────────────────────────────────────────────────
% Usage (once, after % main (standalone preamble for editing)
% vim: nowrap: spell:
% ══════════════════════════════════════════════════════════════════════════════
% Speed Maths --- shared preamble
% Included by every sheet and answer file via:  \input{../../shared/preamble}
% Editing THIS file changes the look of every sheet/answer across every pillar.
% ══════════════════════════════════════════════════════════════════════════════

\usepackage[a4paper, top=25mm, bottom=25mm, left=22mm, right=22mm]{geometry}
\usepackage{amsmath, amssymb}
\usepackage{enumitem}
\usepackage{booktabs}
\usepackage{array}
\usepackage{parskip}
\usepackage{microtype}
\usepackage{fancyhdr}
\usepackage{titlesec}
\usepackage{mdframed}
\usepackage{xcolor}
\usepackage[hidelinks]{hyperref}
\usepackage{graphicx}
\usepackage{tikz}
\usepackage{pgfplots}
\pgfplotsset{compat=1.18}
\usetikzlibrary{calc, angles, quotes}
\tikzset{
  speedgeom/.style={line width=0.8pt, draw=darkgray, line cap=round, line join=round},
  speedaux/.style={line width=0.5pt, dashed, draw=medgray},
  speedpoint/.style={circle, fill=darkgray, inner sep=1.2pt}
}

% ── Colours ───────────────────────────────────────────────────────────────────
\definecolor{darkgray}{gray}{0.15}
\definecolor{medgray}{gray}{0.45}
\definecolor{lightgray}{gray}{0.93}
\definecolor{boxgray}{gray}{0.88}

% ── Section formatting ──────────────────────────────────────────────────────────
\titleformat{\section}[block]
  {\normalfont\large\bfseries}{}
  {0pt}{}[\vspace{2pt}\hrule\vspace{6pt}]
\titlespacing*{\section}{0pt}{14pt}{4pt}

% ── List formatting ─────────────────────────────────────────────────────────────
\setlist[enumerate,1]{label=\textbf{\arabic*.}, leftmargin=2em, itemsep=10pt}

% ── Branding constants (edit here, not per-file) ────────────────────────────────
\newcommand{\SpeedExamLine}{\textbf{Speed Maths:} Competition \& Exam Prep}
\newcommand{\SpeedCredit}{Series by \href{https://www.linkedin.com/in/cerealdev/}{David Akinyele-Aje}}

% ── Header / footer ──────────────────────────────────────────────────────────────
% Usage (once, after \input{../../shared/preamble} and before \begin{document}):
%   \SpeedHeader{Algebra}{3}
% First arg  = pillar name shown as "Speed <Pillar> --- Daily Drill"
% Second arg = worksheet number
\pagestyle{fancy}
\newcommand{\SpeedHeader}[2]{%
  \fancyhf{}%
  \renewcommand{\headrulewidth}{0.4pt}%
  \setlength{\headheight}{14pt}%
  \fancyhead[L]{\small\textbf{Speed #1 --- Daily Drill}}%
  \fancyhead[R]{\small Worksheet \##2}%
  \fancyfoot[L]{\small \SpeedExamLine}%
  \fancyfoot[C]{\small\thepage}%
  \fancyfoot[R]{\small \SpeedCredit}%
  \renewcommand{\footrulewidth}{0.4pt}%
}

% ── Title block (used at the top of every sheet AND answer file) ───────────────
% Usage: \SpeedTitleBlock{Daily Algebra Drill \#3}{Jane Doe}
% %% AGENTS: When generating a new sheet, ask the human contributor for the name they want credited in argument 2.
% For answer files, pass the "--- Answers \& Investigations" suffix in the arg.
\newcommand{\SpeedTitleBlock}[2]{%
  \begin{center}
    {\LARGE\bfseries #1}\\[4pt]
    {\normalsize\itshape Sheet by #2}\\[6pt]
    \hrule\vspace{4pt}
    \begin{tabular}{llcc}
      \toprule
      \textbf{Section} & \textbf{Topic} & \textbf{Questions} & \textbf{Target time}\\
      \midrule
      A & Rapid Recognition          & 10 & 2 min 30 sec \\
      B & Manipulation Drills        & 10 & 8 min        \\
      C & Substitution \& Structure  & 8  & 10 min       \\
      D & Challenge                  & 5  & 15 min       \\
      \bottomrule
    \end{tabular}
    \vspace{4pt}
    \hrule
  \end{center}
}

% ── Sheet metadata: topic + the day's new tools ────────────────────────────────
% Usage (immediately after \SpeedTitleBlock, sheets only):
%   \SpeedMeta{Expanding, factorising, and the identity toolkit}
%             {difference of squares, sum/product form, completing the square}
%
% Arg 1 = the sheet's topic in one line. The website lists this instead of
%         "Sheet 03", so write the thing a reader would actually search for.
% Arg 2 = comma-separated, at most five: the tools this sheet introduces that
%         earlier sheets in the pillar did not. These become the website's tags.
%         Leave out anything the reader already met — the tags are meant to say
%         what is new today, not to summarise the sheet.
%
% Both are parsed straight out of this file by tools/build_website.py, so the
% sheet stays the single source of truth and the site cannot drift from it.
%
% This renders NOTHING. It is a declaration for the build script, not a piece of
% the sheet's layout: adding it to a sheet must not change that sheet's PDF, so
% the 25 existing sheets can be annotated without a single one of them being
% reissued. If the toolkit should also appear on the page, that is a separate
% decision about the sheet design — make it deliberately, here, not by leaving
% this macro printing something nobody asked it to.
\newcommand{\SpeedMeta}[2]{}

% ── Answer / Method / Investigate-further boxes (answer files only) ────────────
\newmdenv[
  backgroundcolor=lightgray,
  linecolor=medgray,
  linewidth=0.5pt,
  leftmargin=0pt, rightmargin=0pt,
  innerleftmargin=8pt, innerrightmargin=8pt,
  innertopmargin=5pt, innerbottommargin=5pt,
  skipabove=4pt, skipbelow=4pt
]{ansbox}

\newmdenv[
  backgroundcolor=white,
  linecolor=darkgray,
  linewidth=0.8pt,
  leftline=true, rightline=false, topline=false, bottomline=false,
  leftmargin=0pt, rightmargin=0pt,
  innerleftmargin=8pt, innerrightmargin=6pt,
  innertopmargin=4pt, innerbottommargin=4pt,
  skipabove=4pt, skipbelow=6pt
]{invbox}

\newcommand{\ans}[1]{%
  \begin{ansbox}
    \textbf{Answer:} #1
  \end{ansbox}}

\newcommand{\method}[1]{%
  {\small\textbf{Method:} #1}\par}

\newcommand{\inv}[1]{%
  \begin{invbox}
    \textbf{\small Investigate further:} {\small #1}
  \end{invbox}}

% ── Misc helpers ────────────────────────────────────────────────────────────────
\newcommand{\qn}[1]{\textbf{#1.}\;}

% Mistake-linking: attach a Top 5 Patterns / Common Traps bullet to the question
% label(s) in this sheet where it actually shows up, e.g. \seealso{B6, D2}
\newcommand{\seealso}[1]{\hfill\textit{\footnotesize(see #1)}}

% ── Graph options macro ─────────────────────────────────────────────────────────
% Usage: \GraphOptions{tikz code for A}{tikz code for B}{tikz code for C}{tikz code for D}
\newcommand{\GraphOptions}[4]{%
  \begin{center}
    \begin{tabular}{cc}
      \textbf{A)} & \textbf{B)} \\
      #1 & #2 \\[10pt]
      \textbf{C)} & \textbf{D)} \\
      #3 & #4
    \end{tabular}
  \end{center}
}

% ── Closing motivational rule + cross-promo footer (sheets only) ────────────────
% Usage: \SpeedClosing{"Quote text."}
\newcommand{\SpeedClosing}[1]{%
  \vspace{20pt}
  \begin{center}
  \rule{0.6\textwidth}{0.4pt}\\[4pt]
  {\small\itshape #1}\\[4pt]
  \rule{0.6\textwidth}{0.4pt}
  \end{center}
  \begin{center}
  {\small Found a mistake or want to contribute a sheet?}\\
  {\small Visit the open-source project at \href{https://github.com/The-CerealDev/speed-maths}{github.com/The-CerealDev/speed-maths}}
  \end{center}
}
 and before \begin{document}):
%   \SpeedHeader{Algebra}{3}
% First arg  = pillar name shown as "Speed <Pillar> --- Daily Drill"
% Second arg = worksheet number
\pagestyle{fancy}
\newcommand{\SpeedHeader}[2]{%
  \fancyhf{}%
  \renewcommand{\headrulewidth}{0.4pt}%
  \setlength{\headheight}{14pt}%
  \fancyhead[L]{\small\textbf{Speed #1 --- Daily Drill}}%
  \fancyhead[R]{\small Worksheet \##2}%
  \fancyfoot[L]{\small \SpeedExamLine}%
  \fancyfoot[C]{\small\thepage}%
  \fancyfoot[R]{\small \SpeedCredit}%
  \renewcommand{\footrulewidth}{0.4pt}%
}

% ── Title block (used at the top of every sheet AND answer file) ───────────────
% Usage: \SpeedTitleBlock{Daily Algebra Drill \#3}{Jane Doe}
% %% AGENTS: When generating a new sheet, ask the human contributor for the name they want credited in argument 2.
% For answer files, pass the "--- Answers \& Investigations" suffix in the arg.
\newcommand{\SpeedTitleBlock}[2]{%
  \begin{center}
    {\LARGE\bfseries #1}\\[4pt]
    {\normalsize\itshape Sheet by #2}\\[6pt]
    \hrule\vspace{4pt}
    \begin{tabular}{llcc}
      \toprule
      \textbf{Section} & \textbf{Topic} & \textbf{Questions} & \textbf{Target time}\\
      \midrule
      A & Rapid Recognition          & 10 & 2 min 30 sec \\
      B & Manipulation Drills        & 10 & 8 min        \\
      C & Substitution \& Structure  & 8  & 10 min       \\
      D & Challenge                  & 5  & 15 min       \\
      \bottomrule
    \end{tabular}
    \vspace{4pt}
    \hrule
  \end{center}
}

% ── Sheet metadata: topic + the day's new tools ────────────────────────────────
% Usage (immediately after \SpeedTitleBlock, sheets only):
%   \SpeedMeta{Expanding, factorising, and the identity toolkit}
%             {difference of squares, sum/product form, completing the square}
%
% Arg 1 = the sheet's topic in one line. The website lists this instead of
%         "Sheet 03", so write the thing a reader would actually search for.
% Arg 2 = comma-separated, at most five: the tools this sheet introduces that
%         earlier sheets in the pillar did not. These become the website's tags.
%         Leave out anything the reader already met — the tags are meant to say
%         what is new today, not to summarise the sheet.
%
% Both are parsed straight out of this file by tools/build_website.py, so the
% sheet stays the single source of truth and the site cannot drift from it.
%
% This renders NOTHING. It is a declaration for the build script, not a piece of
% the sheet's layout: adding it to a sheet must not change that sheet's PDF, so
% the 25 existing sheets can be annotated without a single one of them being
% reissued. If the toolkit should also appear on the page, that is a separate
% decision about the sheet design — make it deliberately, here, not by leaving
% this macro printing something nobody asked it to.
\newcommand{\SpeedMeta}[2]{}

% ── Answer / Method / Investigate-further boxes (answer files only) ────────────
\newmdenv[
  backgroundcolor=lightgray,
  linecolor=medgray,
  linewidth=0.5pt,
  leftmargin=0pt, rightmargin=0pt,
  innerleftmargin=8pt, innerrightmargin=8pt,
  innertopmargin=5pt, innerbottommargin=5pt,
  skipabove=4pt, skipbelow=4pt
]{ansbox}

\newmdenv[
  backgroundcolor=white,
  linecolor=darkgray,
  linewidth=0.8pt,
  leftline=true, rightline=false, topline=false, bottomline=false,
  leftmargin=0pt, rightmargin=0pt,
  innerleftmargin=8pt, innerrightmargin=6pt,
  innertopmargin=4pt, innerbottommargin=4pt,
  skipabove=4pt, skipbelow=6pt
]{invbox}

\newcommand{\ans}[1]{%
  \begin{ansbox}
    \textbf{Answer:} #1
  \end{ansbox}}

\newcommand{\method}[1]{%
  {\small\textbf{Method:} #1}\par}

\newcommand{\inv}[1]{%
  \begin{invbox}
    \textbf{\small Investigate further:} {\small #1}
  \end{invbox}}

% ── Misc helpers ────────────────────────────────────────────────────────────────
\newcommand{\qn}[1]{\textbf{#1.}\;}

% Mistake-linking: attach a Top 5 Patterns / Common Traps bullet to the question
% label(s) in this sheet where it actually shows up, e.g. \seealso{B6, D2}
\newcommand{\seealso}[1]{\hfill\textit{\footnotesize(see #1)}}

% ── Graph options macro ─────────────────────────────────────────────────────────
% Usage: \GraphOptions{tikz code for A}{tikz code for B}{tikz code for C}{tikz code for D}
\newcommand{\GraphOptions}[4]{%
  \begin{center}
    \begin{tabular}{cc}
      \textbf{A)} & \textbf{B)} \\
      #1 & #2 \\[10pt]
      \textbf{C)} & \textbf{D)} \\
      #3 & #4
    \end{tabular}
  \end{center}
}

% ── Closing motivational rule + cross-promo footer (sheets only) ────────────────
% Usage: \SpeedClosing{"Quote text."}
\newcommand{\SpeedClosing}[1]{%
  \vspace{20pt}
  \begin{center}
  \rule{0.6\textwidth}{0.4pt}\\[4pt]
  {\small\itshape #1}\\[4pt]
  \rule{0.6\textwidth}{0.4pt}
  \end{center}
  \begin{center}
  {\small Found a mistake or want to contribute a sheet?}\\
  {\small Visit the open-source project at \href{https://github.com/The-CerealDev/speed-maths}{github.com/The-CerealDev/speed-maths}}
  \end{center}
}
 and before \begin{document}):
%   \SpeedHeader{Algebra}{3}
% First arg  = pillar name shown as "Speed <Pillar> --- Daily Drill"
% Second arg = worksheet number
\pagestyle{fancy}
\newcommand{\SpeedHeader}[2]{%
  \fancyhf{}%
  \renewcommand{\headrulewidth}{0.4pt}%
  \setlength{\headheight}{14pt}%
  \fancyhead[L]{\small\textbf{Speed #1 --- Daily Drill}}%
  \fancyhead[R]{\small Worksheet \##2}%
  \fancyfoot[L]{\small \SpeedExamLine}%
  \fancyfoot[C]{\small\thepage}%
  \fancyfoot[R]{\small \SpeedCredit}%
  \renewcommand{\footrulewidth}{0.4pt}%
}

% ── Title block (used at the top of every sheet AND answer file) ───────────────
% Usage: \SpeedTitleBlock{Daily Algebra Drill \#3}{Jane Doe}
% %% AGENTS: When generating a new sheet, ask the human contributor for the name they want credited in argument 2.
% For answer files, pass the "--- Answers \& Investigations" suffix in the arg.
\newcommand{\SpeedTitleBlock}[2]{%
  \begin{center}
    {\LARGE\bfseries #1}\\[4pt]
    {\normalsize\itshape Sheet by #2}\\[6pt]
    \hrule\vspace{4pt}
    \begin{tabular}{llcc}
      \toprule
      \textbf{Section} & \textbf{Topic} & \textbf{Questions} & \textbf{Target time}\\
      \midrule
      A & Rapid Recognition          & 10 & 2 min 30 sec \\
      B & Manipulation Drills        & 10 & 8 min        \\
      C & Substitution \& Structure  & 8  & 10 min       \\
      D & Challenge                  & 5  & 15 min       \\
      \bottomrule
    \end{tabular}
    \vspace{4pt}
    \hrule
  \end{center}
}

% ── Sheet metadata: topic + the day's new tools ────────────────────────────────
% Usage (immediately after \SpeedTitleBlock, sheets only):
%   \SpeedMeta{Expanding, factorising, and the identity toolkit}
%             {difference of squares, sum/product form, completing the square}
%
% Arg 1 = the sheet's topic in one line. The website lists this instead of
%         "Sheet 03", so write the thing a reader would actually search for.
% Arg 2 = comma-separated, at most five: the tools this sheet introduces that
%         earlier sheets in the pillar did not. These become the website's tags.
%         Leave out anything the reader already met — the tags are meant to say
%         what is new today, not to summarise the sheet.
%
% Both are parsed straight out of this file by tools/build_website.py, so the
% sheet stays the single source of truth and the site cannot drift from it.
%
% This renders NOTHING. It is a declaration for the build script, not a piece of
% the sheet's layout: adding it to a sheet must not change that sheet's PDF, so
% the 25 existing sheets can be annotated without a single one of them being
% reissued. If the toolkit should also appear on the page, that is a separate
% decision about the sheet design — make it deliberately, here, not by leaving
% this macro printing something nobody asked it to.
\newcommand{\SpeedMeta}[2]{}

% ── Answer / Method / Investigate-further boxes (answer files only) ────────────
\newmdenv[
  backgroundcolor=lightgray,
  linecolor=medgray,
  linewidth=0.5pt,
  leftmargin=0pt, rightmargin=0pt,
  innerleftmargin=8pt, innerrightmargin=8pt,
  innertopmargin=5pt, innerbottommargin=5pt,
  skipabove=4pt, skipbelow=4pt
]{ansbox}

\newmdenv[
  backgroundcolor=white,
  linecolor=darkgray,
  linewidth=0.8pt,
  leftline=true, rightline=false, topline=false, bottomline=false,
  leftmargin=0pt, rightmargin=0pt,
  innerleftmargin=8pt, innerrightmargin=6pt,
  innertopmargin=4pt, innerbottommargin=4pt,
  skipabove=4pt, skipbelow=6pt
]{invbox}

\newcommand{\ans}[1]{%
  \begin{ansbox}
    \textbf{Answer:} #1
  \end{ansbox}}

\newcommand{\method}[1]{%
  {\small\textbf{Method:} #1}\par}

\newcommand{\inv}[1]{%
  \begin{invbox}
    \textbf{\small Investigate further:} {\small #1}
  \end{invbox}}

% ── Misc helpers ────────────────────────────────────────────────────────────────
\newcommand{\qn}[1]{\textbf{#1.}\;}

% Mistake-linking: attach a Top 5 Patterns / Common Traps bullet to the question
% label(s) in this sheet where it actually shows up, e.g. \seealso{B6, D2}
\newcommand{\seealso}[1]{\hfill\textit{\footnotesize(see #1)}}

% ── Graph options macro ─────────────────────────────────────────────────────────
% Usage: \GraphOptions{tikz code for A}{tikz code for B}{tikz code for C}{tikz code for D}
\newcommand{\GraphOptions}[4]{%
  \begin{center}
    \begin{tabular}{cc}
      \textbf{A)} & \textbf{B)} \\
      #1 & #2 \\[10pt]
      \textbf{C)} & \textbf{D)} \\
      #3 & #4
    \end{tabular}
  \end{center}
}

% ── Closing motivational rule + cross-promo footer (sheets only) ────────────────
% Usage: \SpeedClosing{"Quote text."}
\newcommand{\SpeedClosing}[1]{%
  \vspace{20pt}
  \begin{center}
  \rule{0.6\textwidth}{0.4pt}\\[4pt]
  {\small\itshape #1}\\[4pt]
  \rule{0.6\textwidth}{0.4pt}
  \end{center}
  \begin{center}
  {\small Found a mistake or want to contribute a sheet?}\\
  {\small Visit the open-source project at \href{https://github.com/The-CerealDev/speed-maths}{github.com/The-CerealDev/speed-maths}}
  \end{center}
}

% Calculus and Graphs only TikZ/pgfplots setup, loaded after shared/preamble.tex
% by every sheet and answer key. Kept out of the shared preamble so this
% pillar's drawing styles cannot change how any other pillar compiles.
\usepackage{needspace}
\usetikzlibrary{calc}
\tikzset{
  cgcurve/.style={line width=0.9pt, draw=darkgray, line cap=round, line join=round},
  cgdash/.style={line width=0.5pt, dashed, draw=medgray},
  cgfill/.style={fill=gray!25, draw=none},
  cghost/.style={line width=0.5pt, dashed, draw=medgray}
}

% ── Section B4 reverse-sketching panels ───────────────────────────────────────
% Every panel is 3.6 cm tall and carries ticks at x = -2, -1, 1, 2, with the
% turning points of the question (-1 and 2) labelled. Every panel has at least
% two flat points, so the answer cannot be found by counting them. Panel D is
% correct: f'(x)=(x+1)(x-2), so f rises, flattens at x=-1, falls, flattens at
% x=2, then rises.
%
% Axes: x in [-2.5,3.35] at 0.72 cm/unit (4.2 cm wide); y in [-1.9,1.95] at
% 0.95 cm/unit (3.66 cm tall). Features are chosen to sit inside that box.
\newcommand{\cgpanelaxes}{%
  \draw[->, line width=0.5pt] (-2.5,0)--(3.35,0);
  \draw[->, line width=0.5pt] (0,-1.9)--(0,1.95);
  \foreach \tx in {-2,-1,1,2}{%
    \draw[line width=0.5pt] (\tx,-0.10)--(\tx,0.10);}%
  \node[below=1.5pt, font=\scriptsize, fill=white, inner sep=0.6pt] at (-1,0) {$-1$};
  \node[below=1.5pt, font=\scriptsize, fill=white, inner sep=0.6pt] at (2,0) {$2$};}
% A: maximum and minimum swapped: minimum at -1, maximum at 2.
\newcommand{\cgBfourA}{%
  \begin{tikzpicture}[x=0.72cm, y=0.95cm]
    \clip (-2.5,-1.9) rectangle (3.35,1.95);
    \cgpanelaxes
    \draw[cgcurve, domain=-2.2:3.2, samples=200]
      plot (\x,{0.55*(-(\x*\x*\x)/3 + (\x*\x)/2 + 2*\x)});
  \end{tikzpicture}}
% B: correct +,-,+ shape but flat points at x=-2 and x=1 instead. The maximum
% sits on the -2 tick and the minimum on the 1 tick.
\newcommand{\cgBfourB}{%
  \begin{tikzpicture}[x=0.72cm, y=0.95cm]
    \clip (-2.5,-1.9) rectangle (3.35,1.95);
    \cgpanelaxes
    \draw[cgcurve, domain=-2.4:2.4, samples=200]
      plot (\x,{0.55*((\x*\x*\x)/3 + (\x*\x)/2 - 2*\x)});
  \end{tikzpicture}}
% C: flat at x=-1 where the curve does NOT turn, then a minimum at x=2.
\newcommand{\cgBfourC}{%
  \begin{tikzpicture}[x=0.72cm, y=0.95cm]
    \clip (-2.5,-1.9) rectangle (3.35,1.95);
    \cgpanelaxes
    \draw[cgcurve, domain=-2.2:3.2, samples=200]
      plot (\x,{0.28*((\x*\x*\x*\x)/4 - 1.5*\x*\x - 2*\x)});
  \end{tikzpicture}}
% D: correct. f'(x)=(x+1)(x-2): maximum at -1, minimum at 2, then rises.
\newcommand{\cgBfourD}{%
  \begin{tikzpicture}[x=0.72cm, y=0.95cm]
    \clip (-2.5,-1.9) rectangle (3.35,1.95);
    \cgpanelaxes
    \draw[cgcurve, domain=-2.2:3.2, samples=200]
      plot (\x,{0.55*((\x*\x*\x)/3 - (\x*\x)/2 - 2*\x)});
  \end{tikzpicture}}
% E: maximum at x=-1, then a flat point at x=2 where the curve does NOT turn
% and keeps falling. f'(x)=-(x+1)(x-2)^2, so f' is negative on both sides of 2.
\newcommand{\cgBfourE}{%
  \begin{tikzpicture}[x=0.72cm, y=0.95cm]
    \clip (-2.5,-1.9) rectangle (3.35,1.95);
    \cgpanelaxes
    \draw[cgcurve, domain=-2.0:2.9, samples=200]
      plot (\x,{0.36*(-(\x*\x*\x*\x)/4 + \x*\x*\x - 4*\x)});
  \end{tikzpicture}}

\SpeedHeader{Calculus and Graphs}{1}

\begin{document}

\SpeedTitleBlock{Daily Calculus \& Graphs Drill \#1 --- Worked Solutions \& Investigations}{Henry Koduthore}
\noindent\textit{New toolkit today: Derivative Notation $\cdot$ The Power Rule $\cdot$ Tangents \& Normals $\cdot$ The Second Derivative.}

%═══════════════════════════════════════════════════════════════════════════════
\section*{Section A \quad Rapid Recognition}
%═══════════════════════════════════════════════════════════════════════════════
\begin{enumerate}

%── A1 ──────────────────────────────────────────────────────────────────────────
\item Find the gradient of the tangent to $y=x^2-3x$ at $x=2$.

\ans{$1$}
\method{$y'=2x-3$, so at $x=2$ the gradient is $4-3=1$.}
\inv{That tangent is $y=x-4$. Solve $x^2-3x=x-4$ to find where the tangent meets the curve again. Why is the intersection a double root, and why can a tangent to a quadratic never cross it at two distinct points?}

%── A2 ──────────────────────────────────────────────────────────────────────────
\item $f(x)=(x+1)(x-3)$. Find $f'(x)$.

\ans{$2x-2$}
\method{Expand first: $(x+1)(x-3)=x^2-2x-3$, so $f'(x)=2x-2$.}
\inv{Expand-first is a habit worth testing. Differentiate $f(x)=(x+1)(x-3)$ after expanding, then try $f(x)=(x+1)(x-3)(x-5)$. How many terms does the expanded cubic have, and which term disappears when you differentiate? What does that tell you about the information $f'$ keeps?}

%── A3 ──────────────────────────────────────────────────────────────────────────
\item $f(x)=\sqrt{x}$. Find $f'(4)$.

\ans{$\frac14$}
\method{Write $f(x)=x^{1/2}$, so $f'(x)=\frac12x^{-1/2}$ and $f'(4)=\frac{1}{2\cdot 2}=\frac14$.}
\inv{Find the secant gradient of $y=\sqrt{x}$ between $0$ and $h$. What happens as $h\to0^+$, and what does that say about the tangent at the origin? Does the secant gradient approach a finite number?}

%── A4 ──────────────────────────────────────────────────────────────────────────
\item $f(x)=\dfrac{1}{x^2}$. Find $f'(x)$.

\ans{$-\frac{2}{x^3}$}
\method{Write $f(x)=x^{-2}$, so $f'(x)=-2x^{-3}=-\frac{2}{x^3}$.}
\inv{State where $f'(x)=-\frac{2}{x^3}$ is positive and where it is negative. Now compare $f(x)=\frac{1}{x}$: there $f'(x)=-\frac{1}{x^2}<0$ everywhere, yet $f(-1)<f(1)$. What does ``decreasing'' mean when the domain is not an interval?}

%── A5 ──────────────────────────────────────────────────────────────────────────
\item $f(x)=x^4-2x^2$. Find $f''(1)$.

\ans{$8$}
\method{$f'(x)=4x^3-4x$, then $f''(x)=12x^2-4$, so $f''(1)=12-4=8$.}
\inv{Here $f''(x)=12x^2-4$, so $f''(0)=-4$ but $f''(2)=44$. What does the sign of $f''$ say about the gradient as $x$ increases? Find where $f''(x)=0$ and describe how the curvature changes either side.}

%── A6 ──────────────────────────────────────────────────────────────────────────
\item $f(x)=x^3-3x^2+4x$. Find $f''(2)$.

\ans{$6$}
\method{$f'(x)=3x^2-6x+4$, so $f''(x)=6x-6$ and $f''(2)=12-6=6$.}
\inv{$f''(x)=6x-6$ changes sign at $x=1$, so the gradient stops falling and starts rising there. Does that make $x=1$ a maximum or minimum of $f$? Evaluate $f'(1)$ and the sign of $f'$ either side before answering.}

%── A7 ──────────────────────────────────────────────────────────────────────────
\item A particle moves along a line with displacement $s=t^3-6t^2+9t$ metres at time $t$ seconds. Find its velocity at $t=2$ and state what the sign tells you.

\ans{$-3$ m/s; moving in the negative direction}
\method{$v=\frac{ds}{dt}=3t^2-12t+9$, so $v(2)=12-24+9=-3$ m/s. The negative sign means the particle is moving in the negative direction.}
\inv{Solve $v=0$: the particle is momentarily at rest at $t=1$ and $t=3$. Evaluate $s$ at both times and the sign of $v$ between them. Describe the motion in words: where does it turn around, and in which direction does it travel on each interval?}

%── A8 ──────────────────────────────────────────────────────────────────────────
\item Find the equation of the normal to $y=x^2$ at $x=1$.

\ans{$y=-\frac{1}{2}x+\frac{3}{2}$}
\method{At $x=1$, $y=1$ and the tangent gradient is $2$; the normal gradient is $-\frac12$. So $y-1=-\frac12(x-1)$, i.e. $y=-\frac12x+\frac32$.}
\inv{The normal at $x=1$ is $y=-\frac12x+\frac32$. Where does it meet the parabola again? Now find every $a$ for which the normal to $y=x^2$ at $x=a$ passes through the origin; how many such normals are there?}

%── A9 ──────────────────────────────────────────────────────────────────────────
\item $f(x)=\dfrac{x^2+4}{\sqrt{x}}$. Find $f'(4)$.

\ans{$\frac{11}{4}$}
\method{Split first: $f(x)=x^{3/2}+4x^{-1/2}$, so $f'(x)=\frac32x^{1/2}-2x^{-3/2}$. At $x=4$: $\frac32\cdot 2-\frac{2}{8}=3-\frac14=\frac{11}{4}$.}
\inv{For $f(x)=\frac{x^2+k}{\sqrt{x}}$, find the value of $k$ for which $f'(4)=0$. Then evaluate $f'$ at $x=3$ and $x=5$ for that $k$: what does the sign pattern say about the gradient near $x=4$?}

%── A10 ─────────────────────────────────────────────────────────────────────────
\item $f(x)=ax^3+bx$, where $f'(1)=5$ and $f''(1)=12$. Find $a$ and $b$.

\ans{$a=2,\ b=-1$}
\method{$f'(x)=3ax^2+b$, so $3a+b=5$. Also $f''(x)=6ax$, so $6a=12$, giving $a=2$ and then $b=5-6=-1$.}
\inv{For $f(x)=ax^3+bx$, show that $f'(-1)=f'(1)$ and $f''(-1)=-f''(1)$ for every $a,b$. So $f'(1)$ and $f''(1)$ can be chosen freely, but $f'(1)$ and $f'(-1)$ cannot. Which pairs of conditions at $x=1$ and $x=-1$ determine $a$ and $b$ uniquely?}

\end{enumerate}

%═══════════════════════════════════════════════════════════════════════════════
\section*{Section B \quad Manipulation Drills}
%═══════════════════════════════════════════════════════════════════════════════
\begin{enumerate}

%── B1 ──────────────────────────────────────────────────────────────────────────
\item Find the equation of the tangent to $y=x^2-4x+5$ at $x=3$.
\begin{itemize}
\item[A)] $y=2x+8$
\item[B)] $y=-\dfrac{1}{2}x+\dfrac{7}{2}$
\item[C)] $y=2x-4$
\item[D)] $y=6x-16$
\item[E)] $y=2x+5$
\end{itemize}

\ans{C) $y=2x-4$}
\method{At $x=3$, $y=9-12+5=2$ and $y'=2x-4$, so the gradient is $2$. Then $y-2=2(x-3)$, i.e.\ $y=2x-4$. Decoys: A comes from the bracket slip $y-2=2(x+3)$; B is the normal; D evaluates $y'$ as $2x$ and forgets the $-4$; E reads the constant term $5$ as the $y$-intercept.}
\inv{Find the point where the tangent to $y=x^2-4x+5$ is perpendicular to $y=2x-4$ ($x=\frac74$). Show that the two tangents meet at a point with $y=\frac34$. Try another pair of perpendicular tangents to this parabola: what do you notice?}

%── B2 ──────────────────────────────────────────────────────────────────────────
\item Find the equation of the normal to $y=x^2-2x$ at $x=3$.
\begin{itemize}
\item[A)] $y=4x-9$
\item[B)] $y=\dfrac{1}{4}x+\dfrac{9}{4}$
\item[C)] $y=-4x+15$
\item[D)] $y=-\dfrac{1}{4}x-\dfrac{9}{4}$
\item[E)] $y=-\dfrac{1}{4}x+\dfrac{15}{4}$
\end{itemize}

\ans{E) $y=-\frac{1}{4}x+\frac{15}{4}$}
\method{At $x=3$, $y=9-6=3$ and the tangent gradient is $2x-2=4$, so the normal gradient is $-\frac14$. Then $y-3=-\frac14(x-3)$, i.e.\ $y=-\frac14x+\frac{15}{4}$. Decoys: A is the tangent; B drops the minus; C negates the gradient without taking the reciprocal; D uses $+3$ instead of $-3$ for the point's $y$-coordinate.}
\inv{Show that the normal meets the parabola again at $x=-\frac54$. For a general $a$, show that the normal at $x=a$ meets $y=x^2-2x$ again at $x=2-a-\frac{1}{2a-2}$. Why must $a=1$ be excluded?}

%── B3 ──────────────────────────────────────────────────────────────────────────
\item For $y=x^2+2x$, what is the average rate of change between $x=1$ and $x=4$?
\begin{itemize}
\item[A)] $21$
\item[B)] $7$
\item[C)] $4$
\item[D)] $10$
\item[E)] $\dfrac{21}{4}$
\end{itemize}

\ans{B) $7$}
\method{$f(1)=3$ and $f(4)=24$, so the average rate is $\frac{24-3}{4-1}=\frac{21}{3}=7$. Decoys: A stops at the total change $21$; C uses the instantaneous rate $f'(1)=4$; D uses $f'(4)=10$; E divides by the upper limit $4$ instead of the interval width $3$.}
\inv{Show that the average rate of change of $y=x^2+2x$ between $x=p$ and $x=q$ is $p+q+2$. Which different intervals give the same average rate?}

%── B4 ──────────────────────────────────────────────────────────────────────────
\item The derivative of $f$ satisfies $f'(x)>0$ for $x<-1$, $f'(-1)=0$, $f'(x)<0$ for $-1<x<2$, $f'(2)=0$, and $f'(x)>0$ for $x>2$. Which sketch could be $y=f(x)$?

\begin{center}
\begin{tabular}{ccc}
\cgBfourA & \cgBfourB & \cgBfourC \\[2pt]
\textbf{A} & \textbf{B} & \textbf{C} \\[10pt]
\cgBfourD & \cgBfourE & \\[2pt]
\textbf{D} & \textbf{E} &
\end{tabular}
\end{center}

\ans{D}
\method{The sign pattern is $+$, $0$, $-$, $0$, $+$: $f$ rises to a high at $x=-1$, falls to a low at $x=2$, then rises. Only panel D does this. Decoys: A swaps the high and the low; B puts the flat points at $-2$ and $1$; C has a flat point at $-1$ where the curve does not turn and only a minimum at $2$; E has a maximum at $-1$, then a flat point at $2$ where the curve does not turn and keeps falling.}
\inv{The correct panel has $f'(x)=(x+1)(x-2)$. Confirm the sign pattern, then explain why a maximum sits where $f'$ changes from $+$ to $-$ and a minimum where it changes from $-$ to $+$. (Day 2 preview: the formal sign chart.)}

%── B5 ──────────────────────────────────────────────────────────────────────────
\item The tangent to $y=x^2$ at $x=a$ passes through $(3,0)$. Find every possible value of $a$.
\begin{itemize}
\item[A)] $a=0$ or $a=6$
\item[B)] $a=6$ only
\item[C)] $a=0$ only
\item[D)] $a=3$
\item[E)] $a=-6$ or $a=0$
\end{itemize}

\ans{A) $a=0$ or $a=6$}
\method{The tangent at $x=a$ is $y=2ax-a^2$. Put $(3,0)$ on it: $0=6a-a^2=a(6-a)$, so $a=0$ or $a=6$. Decoys: B discards $a=0$; C discards $a=6$; D takes the external point's $x$-coordinate as $a$; E comes from the sign slip $0=-6a-a^2$.}
\inv{For a general point $(p,0)$, solve $0=2ap-a^2$. How many tangents to $y=x^2$ pass through $(p,0)$ when $p\neq0$, and what happens when $p=0$?}

%── B6 ──────────────────────────────────────────────────────────────────────────
\item The normal to $y=x^2$ at $x=a$ is parallel to $y=-\dfrac{x}{4}+1$. Find $a$.
\begin{itemize}
\item[A)] $a=-2$
\item[B)] $a=2$
\item[C)] $a=4$
\item[D)] $a=\dfrac{1}{2}$
\item[E)] $a=\dfrac{1}{8}$
\end{itemize}

\ans{B) $a=2$}
\method{The tangent gradient at $x=a$ is $2a$, so the normal gradient is $-\frac{1}{2a}$. Parallel to $y=-\frac{x}{4}+1$ means $-\frac{1}{2a}=-\frac14$, so $2a=4$ and $a=2$. Decoys: A drops the minus; C forgets the factor $2$; D reciprocates the answer; E multiplies by $2$ instead of dividing.}
\inv{Show that the normal to $y=x^2$ at $x=a$ is parallel to the tangent at $x=b$ exactly when $ab=-\frac14$. What does that say about the tangents at $a$ and $b$?}

%── B7 ──────────────────────────────────────────────────────────────────────────
\item The tangent to $y=\sqrt{x}$ at $x=4$ meets the $x$-axis at
\begin{itemize}
\item[A)] $x=4$
\item[B)] $x=1$
\item[C)] $x=-12$
\item[D)] $x=-4$
\item[E)] $x=\dfrac{9}{2}$
\end{itemize}

\ans{D) $x=-4$}
\method{At $x=4$, $y=\sqrt4=2$ and $y'=\frac{1}{2\sqrt{x}}=\frac14$. The tangent is $y-2=\frac14(x-4)$, i.e.\ $y=\frac14x+1$. Setting $y=0$ gives $x=-4$. Decoys: A solves $\frac14x=1$ instead of $\frac14x=-1$; B reports the $y$-intercept; C uses $y(4)=4$ instead of $2$; E uses the normal gradient $-4$ instead of the tangent gradient.}
\inv{Show that the tangent to $y=\sqrt{x}$ at $x=a$ always meets the $x$-axis at $x=-a$. Where does it meet the $y$-axis?}

%── B8 ──────────────────────────────────────────────────────────────────────────
\item For $y=x^3$, find the average rate of change between $x=1$ and $x=1+h$. Simplify your answer, then deduce the instantaneous rate of change at $x=1$.

\ans{average rate $=3+3h+h^2$; instantaneous rate $=3$}
\method{$f(1)=1$ and $f(1+h)=1+3h+3h^2+h^3$. The change is $3h+3h^2+h^3$, so the average rate is $3+3h+h^2$. As $h\to0^+$, this tends to $3$, which is $f'(1)=3$.}
\inv{Repeat from $x=a$ to $x=a+h$. Show that the average rate is $3a^2+3ah+h^2$, then deduce $f'(a)$.}

%── B9 ──────────────────────────────────────────────────────────────────────────
\item Find the equation of the tangent to $y=x^3-3x$ at $x=2$, and find where that tangent meets the curve again.

\ans{$y=9x-16$; again at $(-4,-52)$}
\method{At $x=2$, $y=8-6=2$ and $y'=3x^2-3=9$. The tangent is $y-2=9(x-2)$, i.e.\ $y=9x-16$. Solve $x^3-3x=9x-16$: $x^3-12x+16=(x-2)^2(x+4)=0$, so the other point is $x=-4$, $y=-52$.}
\inv{The tangent at $x=2$ meets the curve again at $x=-4$. Find the tangent at $x=-4$ and check whether it passes through $(2,2)$. What does that show about tangency as a relation between points?}

%── B10 ─────────────────────────────────────────────────────────────────────────
\item The tangent to $y=x^2-2x$ at $P$ is parallel to $y=2x+5$. Find the coordinates of $P$, and the equation of the normal at $P$.

\ans{$P=(2,0)$; normal $y=-\frac{1}{2}x+1$}
\method{$y'=2x-2$. Parallel to $y=2x+5$ means $2x-2=2$, so $x=2$ and $y=4-4=0$. Thus $P=(2,0)$. The normal gradient is $-\frac12$, so $y-0=-\frac12(x-2)$, i.e.\ $y=-\frac12x+1$.}
\inv{Now find the point on $y=x^2-2x$ where the normal is parallel to $y=2x+5$ instead. Solve $-\frac{1}{2x-2}=2$ and interpret the result.}

\end{enumerate}

%═══════════════════════════════════════════════════════════════════════════════
\section*{Section C \quad Substitution \& Structure}
%═══════════════════════════════════════════════════════════════════════════════
\begin{enumerate}

%── C1 ──────────────────────────────────────────────────────────────────────────
\item $f(x)=\dfrac{x^3-8}{x-2}$ for $x\neq 2$. Find $f'(3)$.
\begin{itemize}
\item[A)] $27$
\item[B)] $8$
\item[C)] $4$
\item[D)] $6$
\item[E)] $19$
\end{itemize}

\ans{B) $8$}
\method{Factor the difference of cubes: $x^3-8=(x-2)(x^2+2x+4)$, so $f(x)=x^2+2x+4$ for $x\neq 2$. Then $f'(x)=2x+2$ and $f'(3)=8$. Decoys: A differentiates numerator and denominator separately ($27/1$); C uses the wrong factor $x^2-2x+4$, giving $2x-2$ and $4$; D uses $x^2+4$ (the error $a^3-b^3=(a-b)(a^2+b^2)$), giving $2x$ and $6$; E gives $f(3)=19$.}
\inv{$f(x)$ is the chord gradient of $y=x^3$ between $2$ and $x$. Show that $f(x)\to 12$ as $x\to2$, and explain why that is the derivative of $x^3$ at $2$. Generalise to $\dfrac{x^n-a^n}{x-a}$.}

%── C2 ──────────────────────────────────────────────────────────────────────────
\item The graph of $y=f'(x)$ is a straight line with positive gradient passing through $(1,0)$. Which could be $f(x)$?
\begin{itemize}
\item[A)] $x^2+2x-3$
\item[B)] $-x^2+2x$
\item[C)] $x^3-3x$
\item[D)] $x^2-2x+5$
\item[E)] $2x-2$
\end{itemize}

\ans{D) $x^2-2x+5$}
\method{$f'$ is a straight line $y=m(x-1)$ with $m>0$, so $f$ is a quadratic with $f'(1)=0$ and positive leading coefficient. Only D has $f'=2x-2=2(x-1)$. Decoys: A has $f(1)=0$, confusing a root of $f$ with a root of $f'$ (and $f'(1)=4$); B has $f'=-2(x-1)$, the right line with negative gradient; C has $f'(1)=0$ but $f'=3x^2-3$ is a parabola, not a line; E is $f'$ itself, not $f$.}
\inv{Every $x^2-2x+c$ works. Why can't the graph of $f'$ tell you $c$, and what single extra fact would?}

%── C3 ──────────────────────────────────────────────────────────────────────────
\item A cubic $y=f(x)$ has roots $x=-1$, $2$ and $4$, and $y$-intercept $-16$. Find $f'(0)$.
\begin{itemize}
\item[A)] $-4$
\item[B)] $2$
\item[C)] $4$
\item[D)] $-10$
\item[E)] $-16$
\end{itemize}

\ans{A) $-4$}
\method{Let $f=k(x+1)(x-2)(x-4)$. Then $f(0)=k(1)(-2)(-4)=8k=-16$, so $k=-2$ and $f=-2x^3+10x^2-4x-16$. Hence $f'(x)=-6x^2+20x-4$ and $f'(0)=-4$. Decoys: B assumes $k=1$, giving $f'(0)=2$; C uses the factors $(x-1)(x+2)(x+4)$, which force $k=2$ and $f'(0)=4$; D takes $k$ times the sum of the roots, $-2(5)=-10$; E gives $f(0)=-16$.}
\inv{Show that $f'(-1)=k(-1-2)(-1-4)=-30$ without expanding, by taking the chord gradient $\dfrac{f(x)-f(-1)}{x+1}$ and letting $x\to-1$. Generalise to $f'(p)$ at any root $p$.}

%── C4 ──────────────────────────────────────────────────────────────────────────
\item For $f(x)=x^3-6x^2+9x$, for which $x$ are \textbf{both} $f'(x)<0$ and $f''(x)>0$?
\begin{itemize}
\item[A)] $1<x<3$
\item[B)] $1<x<2$
\item[C)] $x>2$
\item[D)] $x>3$
\item[E)] $2<x<3$
\end{itemize}

\ans{E) $2<x<3$}
\method{$f'=3x^2-12x+9=3(x-1)(x-3)<0$ exactly on $1<x<3$, and $f''=6x-12>0$ exactly for $x>2$. Both hold on $2<x<3$. Decoys: A ignores $f''$; B reverses the sign of $f''$, giving $1<x<2$; C ignores $f'$; D reverses the sign of $f'$, giving $x>3$.}
\inv{On $(2,3)$ the curve is falling but its gradient is increasing. Where is the gradient most negative, and what is that value?}

%── C5 ──────────────────────────────────────────────────────────────────────────
\item The tangents to $y=x^3+kx^2$ at $x=1$ and $x=3$ are parallel. Find $k$.
\begin{itemize}
\item[A)] $6$
\item[B)] $-12$
\item[C)] $-6$
\item[D)] $-\dfrac{13}{4}$
\item[E)] No such $k$ exists
\end{itemize}

\ans{C) $-6$}
\method{$f'=3x^2+2kx$. Parallel tangents at $x=1$ and $x=3$ give $3+2k=27+6k$, so $4k=-24$ and $k=-6$. Decoys: A comes from the sign slip $27-3=6k-2k$; B differentiates $kx^2$ as $kx$, giving $3+k=27+3k$ and $k=-12$; D equates the $y$-values, $1+k=27+9k$, giving $k=-\frac{13}{4}$; E differentiates $kx^2$ as $2k$, so the $k$ terms cancel and no solution exists.}
\inv{Prove that the tangents at $x=p$ and $x=q$ ($p\neq q$) are parallel exactly when $p+q=-\dfrac{2k}{3}$. So parallel tangents pair up symmetrically about $x=-\dfrac{k}{3}$.}

%── C6 ──────────────────────────────────────────────────────────────────────────
\item The line $y=4x+c$ is a tangent to $y=x^3+x$. Find all possible values of $c$.
\begin{itemize}
\item[A)] $c=\pm 1$
\item[B)] $c=\pm 2$
\item[C)] $c=-2$ only
\item[D)] $c=\pm 3$
\item[E)] $c=\pm\dfrac{10\sqrt{3}}{9}$
\end{itemize}

\ans{B) $c=\pm 2$}
\method{$f'=3x^2+1$. A tangent of gradient $4$ needs $3x^2+1=4$, so $x=\pm1$. The contact points are $(1,2)$ and $(-1,-2)$. For $y=4x+c$: $2=4+c$ gives $c=-2$, and $-2=-4+c$ gives $c=2$. Thus $c=\pm2$. Decoys: A reports the $x$-coordinates $\pm1$; C keeps only $x=1$; D takes the contact $y$-values from $y=x^3$ (forgetting the $+x$), giving $\pm1-4(\pm1)=\mp3$; E differentiates $x$ as $0$, so $3x^2=4$, giving $x=\pm\frac{2}{\sqrt3}$ and $c=\mp\frac{10\sqrt3}{9}$.}
\inv{For which gradients $m$ does $y=x^3+x$ have exactly one tangent of gradient $m$? Exactly two? None?}

%── C7 ──────────────────────────────────────────────────────────────────────────
\item The tangent to $y=\dfrac{4}{x}$ at $x=3$ meets the axes at $P$ and $Q$. Find the area of triangle $OPQ$.
\begin{itemize}
\item[A)] $16$
\item[B)] $0$
\item[C)] $4$
\item[D)] $8$
\item[E)] $\dfrac{25}{2}$
\end{itemize}

\ans{D) $8$}
\method{At $x=3$, $y=\frac43$ and $y'=-\frac{4}{x^2}=-\frac49$. The tangent is $y-\frac43=-\frac49(x-3)$, i.e.\ $y=-\frac49x+\frac83$. It meets the axes at $(6,0)$ and $(0,\frac83)$, so the area is $\frac12\cdot6\cdot\frac83=8$. Decoys: A forgets the $\frac12$; B uses $+\frac49$, which sends the tangent through $O$ and gives area $0$; C gives $xy=3\cdot\frac43=4$ at the contact point; E differentiates $\frac4x$ as $-\frac{1}{x^2}$, giving intercepts $15$ and $\frac53$ and area $\frac{25}{2}$.}
\inv{Show that the area is $8$ wherever the point of contact is, and that the contact point is always the midpoint of $PQ$. What is the area for $y=\dfrac{k}{x}$?}

%── C8 ──────────────────────────────────────────────────────────────────────────
\item A cubic $f(x)=ax^3+bx^2+cx+d$ has tangent $y=2x+1$ at $x=1$, with $f''(1)=4$ and $f''(0)=-2$. Find $f(2)$.
\begin{itemize}
\item[A)] $8$
\item[B)] $5$
\item[C)] $6$
\item[D)] $7$
\item[E)] Cannot be determined
\end{itemize}

\ans{A) $8$}
\method{$f''=6ax+2b$. From $f''(0)=-2$, $b=-1$. From $f''(1)=4$, $6a-2=4$, so $a=1$. The tangent at $x=1$ gives both $f'(1)=2$ and $f(1)=3$. Now $f'=3x^2-2x+c$, so $3-2+c=2$ and $c=1$; and $f(1)=1-1+1+d=3$, so $d=2$. Thus $f=x^3-x^2+x+2$ and $f(2)=8-4+2+2=8$. Decoys: B reads $f(2)$ off the tangent ($5$); C ignores $f(1)=3$ and takes $d=0$ ($6$); D writes $f''=6ax+b$, giving $b=-2$, $a=1$, $c=3$, $d=1$ and $f(2)=7$; E treats the tangent as one condition instead of two, leaving $d$ undetermined.}
\inv{The tangent predicts $5$ but the true value is $8$. Show that $f(1+h)-(3+2h)=h^2(2+h)$, and note that the coefficient of $h^2$ is $\frac{f''(1)}{2}$. Is that true for every cubic?}

\end{enumerate}

\needspace{22\baselineskip}
%═══════════════════════════════════════════════════════════════════════════════
\section*{Section D \quad Challenge}
%═══════════════════════════════════════════════════════════════════════════════
\begin{enumerate}

%── D1 ──────────────────────────────────────────────────────────────────────────
\item The chord of $y=x^3$ from $x=-1$ to $x=2$ is parallel to the tangent at $x=c$, where $-1<c<2$. Find all such $c$.
\begin{itemize}
\item[A)] $c=\dfrac12$
\item[B)] $c=-1$ and $c=1$
\item[C)] $c=\sqrt3$
\item[D)] $c=1$ only
\item[E)] $c=\dfrac{2}{\sqrt3}$
\end{itemize}

\ans{D) $c=1$ only}
\method{Chord gradient: $\dfrac{8-(-1)}{2-(-1)}=\dfrac93=3$. A tangent of gradient $3$ needs $3c^2=3$, so $c=\pm1$. The condition $-1<c<2$ is strict, so $c=-1$ is excluded and $c=1$ is the only value. Decoys: A uses the midpoint $\frac12$, which is valid only for quadratics; B ignores the strict inequality; C forgets to divide the rise $9$ by the run $3$, solving $3c^2=9$; E uses the chord from $x=0$, whose gradient is $4$, giving $c=\frac{2}{\sqrt3}$.}
\inv{The value $c=-1$ also gives gradient $3$. Show the chord's line $y=3x+2$ is actually the tangent at $x=-1$, since $x^3-3x-2=(x+1)^2(x-2)$. Prove that for any quadratic the parallel tangent is always at the midpoint, then find where it is for $y=x^3$ on $[0,h]$.}

%── D2 ──────────────────────────────────────────────────────────────────────────
\item The curves $y=x^2$ and $y=-(x-2)^2$ have $y=0$ as a common tangent. Find all their other common tangents.
\begin{itemize}
\item[A)] $y=4x-4$ only
\item[B)] $y=-4x-4$ only
\item[C)] $y=4x-4$ and $y=-4x-4$
\item[D)] $y=2x-1$
\item[E)] There are no others
\end{itemize}

\ans{A) $y=4x-4$ only}
\method{The tangent to $y=x^2$ at $x=a$ is $y=2ax-a^2$. The tangent to $y=-(x-2)^2$ at $x=b$ is $y=-2(b-2)x+b^2-4$. Matching gradients and intercepts gives $a=2-b$ and $-(2-b)^2=b^2-4$, so $2b^2-4b=0$ and $b=0$ or $2$. The root $b=2$ gives $a=0$ and the common tangent $y=0$; the other root $b=0$ gives $a=2$ and $y=4x-4$. Decoys: B takes the second gradient as $+2(b-2)$, giving $y=-4x-4$; C assumes mirror symmetry and keeps both; D forces both points of contact to have the same $x$-coordinate, $2a=-2(a-2)$, giving $a=1$ and $y=2x-1$; E writes the first tangent as $y=2ax+a^2$, which leaves only $y=0$.}
\inv{For $y=x^2$ and $y=-(x-k)^2$, show the common tangents are $y=0$ and $y=2kx-k^2$, touching at $x=k$ and $x=0$ respectively. What happens as $k\to0$? What if the second curve is shifted to $y=-(x-k)^2+m$?}

%── D3 ──────────────────────────────────────────────────────────────────────────
\item The curve $y=f(x)=x^3+bx^2+cx+d$ has horizontal tangents at $x=-1$ and $x=3$. Find $f(-1)-f(3)$.
\begin{itemize}
\item[A)] $-8$
\item[B)] $64$
\item[C)] $-32$
\item[D)] Cannot be determined without $d$
\item[E)] $32$
\end{itemize}

\ans{E) $32$}
\method{Write $f'=3x^2+2bx+c$. Horizontal tangents at $x=-1$ and $x=3$ give $f'=3(x+1)(x-3)=3x^2-6x-9$, so $b=-3$, $c=-9$. Then $f(-1)=5+d$ and $f(3)=-27+d$, so $f(-1)-f(3)=32$ for every $d$. Decoys: A forgets the factor $3$, using $f'=(x+1)(x-3)$, so $b=-1$, $c=-3$; with the cubic still monic this gives $f(-1)-f(3)=-8$; B assumes the gap is $(3-(-1))^3=64$; C computes $f(3)-f(-1)=-32$; D does not notice that $d$ cancels.}
\inv{Show that for $y=x^3+bx^2+cx+d$ with horizontal tangents a distance $D$ apart, the vertical gap is always $\dfrac{D^3}{2}$. Do NOT use integration (that is Day 4): translate so the horizontal tangents are at $x=\pm h$, giving $y=x^3-3h^2x+\text{const}$. What changes if the leading coefficient is $k$?}

%── D4 ──────────────────────────────────────────────────────────────────────────
\item Several tangents to $y=x^3-3x$ pass through the point $(2,-2)$. Find the sum of the $x$-coordinates of all their points of contact.
\begin{itemize}
\item[A)] $1$
\item[B)] $3$
\item[C)] $2$
\item[D)] $2+\sqrt3$
\item[E)] $-3$
\end{itemize}

\ans{B) $3$}
\method{The tangent at $x=a$ is $y=(3a^2-3)x-2a^3$. Passing through $(2,-2)$ gives $-2=2(3a^2-3)-2a^3$, so $a^3-3a^2+2=0$. Since $a=1$ is a root, $(a-1)(a^2-2a-2)=0$, giving $a=1$ and $a=1\pm\sqrt3$. The sum of the three contact $x$-coordinates is $3$. Decoys: A finds only the horizontal tangent $a=1$; C keeps only the two roots of the quadratic factor, whose sum is $2$; D discards the negative root $1-\sqrt3$ but keeps $a=1$ and $1+\sqrt3$, giving $2+\sqrt3$; E makes a sign error in the sum of the roots, giving $-3$.}
\inv{$(2,2)$ lies on the curve. Show the equation for $a$ becomes $(a+1)(a-2)^2=0$. How many distinct tangents pass through $(2,2)$, and why is $a=2$ a double root? (Day 2 preview.) For which points are there exactly three tangents?}

%── D5 ──────────────────────────────────────────────────────────────────────────
\item Let $a\neq0$.
\begin{enumerate}
\item[(a)] Show that the tangent to $y=x^3$ at $x=a$ meets the curve again at $x=-2a$.
\item[(b)] Show that the gradient of the curve at $x=-2a$ is four times the gradient at $x=a$.
\item[(c)] The tangent at $x=-2a$ meets the curve again at $x=b$. Find $b$, and the gradient there in terms of $a$.
\item[(d)] Explain why the condition $a\neq0$ is needed.
\end{enumerate}

\ans{$b=4a$; gradient $48a^2$}
\method{(a) The tangent at $x=a$ is $y=3a^2x-2a^3$. Subtracting from $y=x^3$ gives $x^3-3a^2x+2a^3$. Factor and expand to verify: $(x-a)^2(x+2a)=(x^2-2ax+a^2)(x+2a)=x^3-3a^2x+2a^3$. Since $a\neq0$, the factor $x+2a$ gives a second intersection at $x=-2a$. (b) The gradient at $x=-2a$ is $3(-2a)^2=12a^2=4\cdot3a^2$, four times the gradient at $x=a$. (c) Replacing $a$ by $-2a$ in part (a), the tangent at $x=-2a$ meets the curve again at $x=-2(-2a)=4a$, so $b=4a$; the gradient there is $3(4a)^2=48a^2$. (d) If $a=0$, the tangent is $y=0$, and $x^3=0$ has only the triple root $x=0$; there is no second point.}
\inv{For $y=x^3+px^2+qx+r$, show the tangent at $x=a$ meets the curve again at $x=-2a-p$. For which $a$ does it not meet the curve again, and what is special about that point?}

\end{enumerate}

%═══════════════════════════════════════════════════════════════════════════════
\section*{Top 5 Patterns --- Worksheet \#1}
\begin{enumerate}[leftmargin=2em,itemsep=4pt]
  \item \textbf{Rewrite as powers before differentiating:} $\sqrt{x}=x^{1/2}$, $1/x^2=x^{-2}$; split or factorise fractions first. \seealso{A3, A4, A9, B7, C1, C7}
  \item \textbf{Tangent at $x=a$:} $y=f'(a)(x-a)+f(a)$. Normal gradient: $-1/f'(a)$. \seealso{A1, A8, B1, B2, B6, B10}
  \item \textbf{Tangency means a repeated root:} substitute the line into the curve, factor out $(x-a)^2$, and the remaining factor gives where the tangent meets the curve again. \seealso{B9, D1, D5}
  \item \textbf{Unknown point of contact:} write the tangent at a general point $x=a$, impose the condition (passes through a point, has a given gradient, is shared by two curves), then solve for $a$, keeping every root. \seealso{B5, C6, D2, D4}
  \item \textbf{Each condition is one equation:} a value of $f$, $f'$ or $f''$ gives one equation, and a tangent line gives two (value and gradient). Count conditions against unknowns. \seealso{A10, C3, C8, D3}
\end{enumerate}

\section*{Common Traps --- Worksheet \#1}
\begin{enumerate}[leftmargin=2em,itemsep=4pt]
  \item \textbf{Differentiating a quotient top and bottom separately:} simplify first. \seealso{A9, C1}
  \item \textbf{Normal gradient:} you must negate AND take the reciprocal; doing only one gives the wrong line. \seealso{B2, B6}
  \item \textbf{Losing roots:} keep obvious roots ($a=0$ in B5, $a=1$ in D4) and respect strict inequalities ($c=-1$ is excluded in D1). \seealso{B5, D1, D4}
  \item \textbf{Building $f'$ from turning points:} $f'(x)=3(x-p)(x-q)$ for a cubic with leading coefficient $1$; forgetting the $3$ changes every coefficient. \seealso{C5, D3}
  \item \textbf{A tangent line is two conditions, not one:} it fixes both $f(a)$ and $f'(a)$. \seealso{C8}
\end{enumerate}

\SpeedClosing{``Simplify first, then differentiate --- never the other way round.''}

\end{document}
